\documentclass[11pt, a4paper]{article}

\usepackage[utf8]{inputenc}
\usepackage[T1]{fontenc}
\usepackage{amsmath, amssymb, amsthm}
\usepackage{mathtools}
\usepackage{geometry}
\usepackage{enumitem}
\usepackage{xcolor}
\usepackage{tcolorbox}
\usepackage{hyperref}

\geometry{margin=1in}

\newtheorem{proposition}{Proposition}
\newtheorem{remark}{Remark}
\theoremstyle{definition}
\newtheorem{definition}{Definition}
\newtheorem{goal}{Goal}

\tcbuselibrary{breakable, skins}

\newtcolorbox{keyidea}[1][]{
  colback=blue!5!white,
  colframe=blue!50!black,
  title={Key Idea},
  fonttitle=\bfseries,
  breakable,
  #1
}

\newtcolorbox{calculation}[1][]{
  colback=green!3!white,
  colframe=green!40!black,
  title={Detailed Calculation},
  fonttitle=\bfseries,
  breakable,
  #1
}

\newtcolorbox{intuition}[1][]{
  colback=orange!5!white,
  colframe=orange!50!black,
  title={Intuition},
  fonttitle=\bfseries,
  breakable,
  #1
}

\title{\textbf{Mini-Lecture: The Gaussian AWGN Case Study} \\[6pt]
\large A Step-by-Step Walkthrough of Section 5 of \\
``Generative Modeling from Black-box Corruptions \\via Self-Consistent Stochastic Interpolants''}
\author{Prepared as a pedagogical companion}
\date{}

\begin{document}
\maketitle

\tableofcontents
\newpage

%=============================================================================
\section{What This Section Is About}
%=============================================================================

The paper introduces a method called \textbf{Self-Consistent Stochastic Interpolants (SCSI)} for learning a generative model of clean data $\pi$ when you only have access to corrupted observations $\mu$ and a black-box forward corruption model $\mathcal{F}$.

Section 5 takes the \emph{simplest possible case} --- a Gaussian prior corrupted by additive white Gaussian noise (AWGN) --- and works out \emph{everything} in closed form. This lets us:
\begin{enumerate}[label=(\roman*)]
    \item Verify that the iterative scheme actually converges.
    \item Compute the exact convergence rate.
    \item Compare SCSI (with ODE transport) to the classical EM algorithm.
    \item Check the theoretical assumptions from Section 4.
\end{enumerate}

%=============================================================================
\section{Setup}
%=============================================================================

\subsection{The True Data Distribution}

We assume the clean data follows a Gaussian distribution:
\[
    \pi = \mathcal{N}(b, \Sigma), \qquad \Sigma \succ 0 \;\text{(positive definite)}.
\]

\subsection{The Forward (Corruption) Channel}

The AWGN channel adds independent Gaussian noise:
\[
    y = \mathcal{F}(x) = x + w, \qquad w \sim \mathcal{N}(0, I).
\]

\begin{intuition}
We've \emph{normalized} the noise variance to $I$ without loss of generality. If the original channel were $y = x + \sigma \xi$ with $\xi \sim \mathcal{N}(0,I)$, we could rescale $x \to x/\sigma$ to get unit noise. The ``signal-to-noise ratio'' (SNR) is then encoded in $\Sigma$ --- eigenvalues of $\Sigma$ that are large relative to 1 correspond to high SNR directions, and small eigenvalues correspond to low SNR.
\end{intuition}

\subsection{The Observation Distribution}

Since $y = x + w$ with $x \sim \mathcal{N}(b, \Sigma)$ and $w \sim \mathcal{N}(0, I)$ independent:
\[
    \mu = \text{Law}(y) = \mathcal{N}(b, \Sigma + I).
\]
This is just the convolution $\mu = \pi \star \gamma_1$, where $\gamma_1 = \mathcal{N}(0, I)$.

%=============================================================================
\section{The Stochastic Interpolant for This Problem}
%=============================================================================

\subsection{Choice of Interpolant Schedule}

The paper uses the specific interpolant:
\[
    I_t = (1 - \sqrt{t})\, x + \sqrt{t}\, \mathcal{F}(x) = (1-\sqrt{t})\,x + \sqrt{t}\,(x + w) = x + \sqrt{t}\, w.
\]

\begin{calculation}[title={Simplifying the interpolant}]
Let's verify this step by step. We have $\alpha_t = 1-\sqrt{t}$ and $\beta_t = \sqrt{t}$, and we set $\gamma_t = 0$ (no extra noise term). Then:
\begin{align}
    I_t &= \alpha_t \, x + \beta_t \, \mathcal{F}(x) \notag \\
        &= (1 - \sqrt{t})\, x + \sqrt{t}\,(x + w) \notag \\
        &= x - \sqrt{t}\, x + \sqrt{t}\, x + \sqrt{t}\, w \notag \\
        &= x + \sqrt{t}\, w. \label{eq:It}
\end{align}
This is a beautiful simplification: the interpolant is just the clean signal $x$ plus a \emph{scaled} version of the noise $w$.
\end{calculation}

\subsection{The Law of $I_t$}

\begin{keyidea}
For a \emph{general} Gaussian measure $\tilde{\pi} = \mathcal{N}(\tilde{b}, \tilde{\Sigma})$ replacing $\pi$, the law of $I_t$ is:
\[
    \tilde{\pi}_t = \text{Law}(I_t) = \mathcal{N}(\tilde{b}, \, \tilde{\Sigma} + t\, I).
\]
\end{keyidea}

\begin{calculation}[title={Deriving the law of $I_t$}]
If $x \sim \mathcal{N}(\tilde{b}, \tilde{\Sigma})$ and $w \sim \mathcal{N}(0, I)$ are independent, then $I_t = x + \sqrt{t}\, w$ is a sum of independent Gaussians:
\begin{align}
    \mathbb{E}[I_t] &= \mathbb{E}[x] + \sqrt{t}\,\mathbb{E}[w] = \tilde{b} + 0 = \tilde{b}, \notag \\
    \text{Cov}(I_t) &= \text{Cov}(x) + t\,\text{Cov}(w) = \tilde{\Sigma} + t\, I. \notag
\end{align}
So indeed $I_t \sim \mathcal{N}(\tilde{b},\, \tilde{\Sigma} + tI)$.

\medskip\noindent\textbf{Boundary conditions:}
\begin{itemize}
    \item At $t=0$: $I_0 = x \sim \mathcal{N}(\tilde{b}, \tilde{\Sigma})$ \;$\checkmark$\; (the ``clean'' distribution).
    \item At $t=1$: $I_1 = x + w \sim \mathcal{N}(\tilde{b}, \tilde{\Sigma} + I)$ \;$\checkmark$\; (the ``corrupted'' distribution $\mathcal{K}\tilde{\pi}$).
\end{itemize}
\end{calculation}

\subsection{Connection to the Heat Equation}

The evolution $\tilde{\pi}_t = \mathcal{N}(\tilde{b}, \tilde{\Sigma} + tI)$ corresponds exactly to the \textbf{heat equation}:
\[
    \partial_t \tilde{\pi}_t = \tfrac{1}{2}\Delta \tilde{\pi}_t, \qquad \tilde{\pi}_0 = \tilde{\pi}.
\]

\begin{intuition}
Running the heat equation for time $t$ on a Gaussian $\mathcal{N}(\tilde{b}, \tilde{\Sigma})$ simply adds $tI$ to the covariance --- this is exactly what Gaussian convolution does. The forward corruption process ``is'' the heat equation from $t=0$ to $t=1$.
\end{intuition}

%=============================================================================
\section{The Score Function}
%=============================================================================

For a Gaussian $\tilde{\pi}_t = \mathcal{N}(\tilde{b}, \tilde{\Sigma} + tI)$, the score (gradient of the log-density) is:
\[
    \nabla \log \tilde{\pi}_t(x) = -(\tilde{\Sigma} + tI)^{-1}(x - \tilde{b}).
\]

\begin{calculation}[title={Deriving the Gaussian score}]
The density of $\mathcal{N}(\tilde{b}, C)$ with $C = \tilde{\Sigma}+tI$ is:
\[
    \tilde{\pi}_t(x) = \frac{1}{(2\pi)^{d/2}\det(C)^{1/2}} \exp\!\Bigl(-\tfrac{1}{2}(x-\tilde{b})^\top C^{-1}(x-\tilde{b})\Bigr).
\]
Taking the log:
\[
    \log \tilde{\pi}_t(x) = \text{const} - \tfrac{1}{2}(x-\tilde{b})^\top C^{-1}(x-\tilde{b}).
\]
Taking the gradient with respect to $x$:
\[
    \nabla_x \log \tilde{\pi}_t(x) = -C^{-1}(x-\tilde{b}) = -(\tilde{\Sigma}+tI)^{-1}(x-\tilde{b}).
\]
This is an \textbf{affine} function of $x$, which will be crucial for keeping everything in closed form.
\end{calculation}

%=============================================================================
\section{The Reverse Transport SDE}
%=============================================================================

\subsection{General Form}

The paper uses a reverse-time SDE to transport observations (at $t=1$) back to clean data (at $t=0$). In the notation of the paper, switching to the reverse time variable $\tilde{t}$ (where $\tilde{t}=0$ is the observation end and $\tilde{t}=1$ is the data end):
\[
    dY_{\tilde{t}} = \frac{1+\epsilon}{2}\,\nabla\log\tilde{\pi}_{1-\tilde{t}}(Y_{\tilde{t}})\,d\tilde{t} + \sqrt{\epsilon}\,dW_{\tilde{t}}.
\]

\begin{calculation}[title={Where does this SDE come from?}]
The general reverse-time SDE from Section 2.1 (Eq.~2 in the paper) has drift $b + \epsilon \gamma_t^{-1} g$, where $b$ is the velocity field and $g$ is the denoiser. For the heat equation channel, the connection between score and these fields gives a total drift of $\frac{1+\epsilon}{2}\nabla\log\tilde{\pi}_{1-\tilde{t}}$.

\medskip
\textbf{The key point:} the reverse process is an Ornstein--Uhlenbeck (OU) process because $\nabla\log\tilde{\pi}_{1-\tilde{t}}$ is \emph{affine} in $x$ (see previous section). OU processes can be solved exactly.

\medskip
\textbf{Special case $\epsilon=0$ (ODE):}
\[
    dY_{\tilde{t}} = \tfrac{1}{2}\nabla\log\tilde{\pi}_{1-\tilde{t}}(Y_{\tilde{t}})\,d\tilde{t}.
\]
This is a deterministic ODE --- no randomness, just follow the score.
\end{calculation}

\subsection{Substituting the Gaussian Score}

Substituting the score $\nabla\log\tilde{\pi}_{1-\tilde{t}}(x) = -(\tilde{\Sigma}+(1-\tilde{t})I)^{-1}(x - \tilde{b})$:
\[
    dY_{\tilde{t}} = -\frac{1+\epsilon}{2}(\tilde{\Sigma}+(1-\tilde{t})I)^{-1}(Y_{\tilde{t}} - \tilde{b})\,d\tilde{t} + \sqrt{\epsilon}\,dW_{\tilde{t}}.
\]
This is a \textbf{time-varying linear SDE} (Ornstein--Uhlenbeck type) and can be solved in closed form.

%=============================================================================
\section{Solving the Reverse SDE: The Solution Map}
%=============================================================================

\subsection{The Fundamental Solution Map $\Phi(t,s)$}

The paper defines the solution (or ``propagator'') map:
\[
    \boxed{\;\Phi(t,s) = \bigl(\tilde{\Sigma}+(1-t)I\bigr)^{(1+\epsilon)/2} \bigl(\tilde{\Sigma}+(1-s)I\bigr)^{-(1+\epsilon)/2}\;}
\]
which maps the state from time $s$ to time $t$.

\begin{calculation}[title={Verifying that $\Phi$ solves the homogeneous equation}]
Consider the homogeneous (noise-free) part of the SDE. We need to solve:
\[
    \frac{d}{d\tilde{t}}\,\Phi({\tilde{t}},s) = -\frac{1+\epsilon}{2}(\tilde{\Sigma}+(1-\tilde{t})I)^{-1}\,\Phi({\tilde{t}},s), \qquad \Phi(s,s) = I.
\]
Since all the matrices $\tilde{\Sigma}+(1-\tilde{t})I$ commute with each other for different values of $\tilde{t}$ (they are all functions of the \emph{same} matrix $\tilde{\Sigma}$), we can use the matrix exponential approach.

\medskip
\textbf{Key trick:} Define $A(\tilde{t}) = \tilde{\Sigma}+(1-\tilde{t})I$. Note that:
\[
    \frac{d}{d\tilde{t}}\, A(\tilde{t})^{(1+\epsilon)/2} = \frac{1+\epsilon}{2}\,A(\tilde{t})^{(1+\epsilon)/2-1}\cdot(+I) = \frac{1+\epsilon}{2}\,A(\tilde{t})^{(1+\epsilon)/2}\, A(\tilde{t})^{-1},
\]
where we used the chain rule and $\frac{d}{d\tilde{t}}A(\tilde{t}) = -(-I) = +I$...

\medskip
Wait, let's be more careful. We have $A(\tilde{t}) = \tilde{\Sigma} + (1-\tilde{t})I$, so:
\[
    \frac{dA}{d\tilde{t}} = -I.
\]
Now, for a matrix function $f(A) = A^p$ where $A$ depends on a scalar parameter and all values of $A$ commute:
\[
    \frac{d}{d\tilde{t}} A(\tilde{t})^p = p\, A(\tilde{t})^{p-1}\frac{dA}{d\tilde{t}} = p\,A(\tilde{t})^{p-1}(-I) = -p\,A(\tilde{t})^{p-1}.
\]
So with $p = (1+\epsilon)/2$:
\[
    \frac{d}{d\tilde{t}} A(\tilde{t})^{(1+\epsilon)/2} = -\frac{1+\epsilon}{2}\,A(\tilde{t})^{(1+\epsilon)/2 - 1}.
\]
Now compute:
\begin{align}
    \frac{d}{d\tilde{t}}\Phi(\tilde{t},s) &= \frac{d}{d\tilde{t}}\Bigl[A(\tilde{t})^{(1+\epsilon)/2}\Bigr]\cdot A(s)^{-(1+\epsilon)/2} \notag \\
    &= -\frac{1+\epsilon}{2}\,A(\tilde{t})^{(1+\epsilon)/2-1}\cdot A(s)^{-(1+\epsilon)/2} \notag \\
    &= -\frac{1+\epsilon}{2}\,A(\tilde{t})^{-1}\cdot \underbrace{A(\tilde{t})^{(1+\epsilon)/2}\cdot A(s)^{-(1+\epsilon)/2}}_{\Phi(\tilde{t},s)} \notag \\
    &= -\frac{1+\epsilon}{2}\,A(\tilde{t})^{-1}\,\Phi(\tilde{t},s). \notag
\end{align}
This is exactly the ODE $\dot{\Phi} = -\frac{1+\epsilon}{2}(\tilde{\Sigma}+(1-\tilde{t})I)^{-1}\Phi$. Also $\Phi(s,s) = A(s)^{(1+\epsilon)/2}A(s)^{-(1+\epsilon)/2} = I$. $\checkmark$
\end{calculation}

\subsection{The Full Solution via Variation of Constants}

Using the solution map, the full SDE solution from $\tilde{t}=0$ (observations) to $\tilde{t}=1$ (data) is:
\begin{equation}\label{eq:Y1}
    Y_1 = \tilde{b} + \Phi(1,0)(Y_0 - \tilde{b}) + \sqrt{\epsilon}\int_0^1 \Phi(1,s)\,dW_s.
\end{equation}

\begin{calculation}[title={Deriving Eq.~\eqref{eq:Y1} by variation of constants}]
Write $Z_{\tilde{t}} = Y_{\tilde{t}} - \tilde{b}$ (center the process). Then $Z$ satisfies:
\[
    dZ_{\tilde{t}} = -\frac{1+\epsilon}{2}A(\tilde{t})^{-1}Z_{\tilde{t}}\,d\tilde{t} + \sqrt{\epsilon}\,dW_{\tilde{t}}.
\]
The variation-of-constants formula for a linear SDE $dZ = F(\tilde{t})Z\,d\tilde{t} + \sigma\,dW$ gives:
\[
    Z_{\tilde{t}} = \Phi(\tilde{t},0)\,Z_0 + \sqrt{\epsilon}\int_0^{\tilde{t}} \Phi(\tilde{t},s)\,dW_s.
\]
At $\tilde{t}=1$:
\[
    Z_1 = \Phi(1,0)\,Z_0 + \sqrt{\epsilon}\int_0^1\Phi(1,s)\,dW_s.
\]
Substituting back $Y = Z + \tilde{b}$:
\[
    Y_1 = \tilde{b} + \Phi(1,0)(Y_0 - \tilde{b}) + \sqrt{\epsilon}\int_0^1\Phi(1,s)\,dW_s.
\]
\end{calculation}

\begin{intuition}
Formula \eqref{eq:Y1} has three parts:
\begin{enumerate}
    \item $\tilde{b}$: the mean of the current prior guess.
    \item $\Phi(1,0)(Y_0 - \tilde{b})$: a ``shrunken'' version of how far the observation deviates from the mean. The map $\Phi(1,0)$ is a contraction that pulls $Y_0$ toward $\tilde{b}$.
    \item The stochastic integral: noise injected during the reverse diffusion (absent when $\epsilon = 0$).
\end{enumerate}
\end{intuition}

%=============================================================================
\section{The Iterative Update Rules}
%=============================================================================

\subsection{How the Iteration Works}

\begin{goal}
We start with an initial guess $\pi_0 = \mathcal{N}(b_0, \Sigma_0)$ and want to iteratively converge to $\pi = \mathcal{N}(b, \Sigma)$.
\end{goal}

At iteration $k$, our current guess is $\pi_k = \mathcal{N}(b_k, \Sigma_k)$.

\medskip\noindent\textbf{One iteration of SCSI:}
\begin{enumerate}
    \item Take observation $Y_0 \sim \mu = \mathcal{N}(b, \Sigma+I)$ (these are the real corrupted data).
    \item Run the reverse SDE \eqref{eq:Y1} with $\tilde{b} = b_k$, $\tilde{\Sigma} = \Sigma_k$ to get $Y_1 = \pi_{k+1}$.
\end{enumerate}

\subsection{Computing $\pi_{k+1}$}

Since $Y_0 \sim \mathcal{N}(b, \Sigma + I)$ and the transformation \eqref{eq:Y1} is affine, $Y_1$ is also Gaussian, so $\pi_{k+1} = \mathcal{N}(b_{k+1}, \Sigma_{k+1})$.

\subsubsection{The Matrix $B_k$}

Define:
\[
    \boxed{\;B_k = \bigl(\Sigma_k(\Sigma_k + I)^{-1}\bigr)^{(1+\epsilon)/2}\;}
\]

\begin{calculation}[title={Where does $B_k$ come from?}]
Recall $\Phi(1,0) = A(1)^{(1+\epsilon)/2}\,A(0)^{-(1+\epsilon)/2}$ with $A(\tilde{t}) = \Sigma_k + (1-\tilde{t})I$.

\begin{itemize}
    \item $A(0) = \Sigma_k + I$
    \item $A(1) = \Sigma_k + 0 \cdot I = \Sigma_k$
\end{itemize}

Therefore:
\begin{align}
    \Phi(1,0) &= \Sigma_k^{(1+\epsilon)/2}\,(\Sigma_k + I)^{-(1+\epsilon)/2} \notag \\
              &= \bigl(\Sigma_k(\Sigma_k+I)^{-1}\bigr)^{(1+\epsilon)/2} \notag \\
              &= B_k. \notag
\end{align}
The last equality uses the fact that $\Sigma_k$ and $\Sigma_k + I$ commute (they are simultaneously diagonalizable), so we can combine the exponents.

\medskip
\textbf{Observation:} Since $0 \prec \Sigma_k(\Sigma_k+I)^{-1} \prec I$ (all eigenvalues are between 0 and 1), we have $0 \prec B_k \prec I$ for any $\epsilon \geq 0$.
\end{calculation}

\subsubsection{Mean Update}

\begin{calculation}[title={Deriving $b_{k+1}$}]
From \eqref{eq:Y1}:
\begin{align}
    b_{k+1} &= \mathbb{E}[Y_1] = b_k + B_k(\mathbb{E}[Y_0] - b_k) + 0 \notag \\
             &= b_k + B_k(b - b_k). \label{eq:mean_update}
\end{align}
Here we used:
\begin{itemize}
    \item $\mathbb{E}[Y_0] = b$ (the true mean of the observations $\mu = \mathcal{N}(b, \Sigma+I)$; note the mean of $\mu$ is $b$, not $b_k$).
    \item The stochastic integral has zero mean.
\end{itemize}

We can rewrite this as:
\[
    b_{k+1} - b = (I - B_k)(b_k - b),
\]
or equivalently, defining $r_k = b - b_k$ (the mean error):
\[
    \boxed{\;r_{k+1} = (I - B_k)\,r_k.\;}
\]
Since $0 \prec B_k \prec I$, we have $0 \prec I - B_k \prec I$, which means $\|r_{k+1}\| < \|r_k\|$ --- the mean converges!
\end{calculation}

\subsubsection{Covariance Update}

\begin{calculation}[title={Deriving $\Sigma_{k+1}$}]
From \eqref{eq:Y1}, $Y_1 - b_{k+1}$ equals:
\[
    B_k(Y_0 - b) + \sqrt{\epsilon}\int_0^1 \Phi(1,s)\,dW_s.
\]
So the covariance is (since $Y_0$ and $W$ are independent):
\begin{align}
    \Sigma_{k+1} &= B_k\,\text{Cov}(Y_0)\,B_k + \epsilon\int_0^1 \Phi(1,s)\Phi(1,s)^\top\,ds \notag \\
                 &= B_k\,(\Sigma + I)\,B_k + \epsilon \int_0^1 \Phi(1,s)^2\,ds. \label{eq:cov_intermediate}
\end{align}
(Note: all matrices commute, so $\Phi(1,s)^\top = \Phi(1,s)$.)

\medskip
Now, one can show that the stochastic integral term combines with the first term to give the clean formula. Indeed, the paper states the result as:
\[
    \boxed{\;\Sigma_{k+1} = \Sigma_k + B_k(\Sigma - \Sigma_k)B_k.\;}
\]
Let's verify this is consistent. Expanding:
\begin{align}
    \Sigma_k + B_k(\Sigma - \Sigma_k)B_k &= \Sigma_k + B_k\Sigma B_k - B_k\Sigma_k B_k \notag \\
    &= (I - B_k^2)\Sigma_k + B_k\Sigma B_k. \notag
\end{align}

Now, from \eqref{eq:cov_intermediate}, we need:
\[
    B_k(\Sigma + I)B_k + \text{(noise term)} = (I - B_k^2)\Sigma_k + B_k \Sigma B_k.
\]
The noise term must equal:
\[
    (I - B_k^2)\Sigma_k + B_k\Sigma B_k - B_k(\Sigma+I)B_k = (I - B_k^2)\Sigma_k - B_k^2.
\]
Using $B_k^2 = (\Sigma_k(\Sigma_k+I)^{-1})^{1+\epsilon}$, and recalling that $\Sigma_k$ and $B_k$ commute, one can verify this identity holds (the It\^o integral contribution exactly fills the gap). The full calculation is somewhat involved, so the paper takes this as the result.

\medskip
Defining the covariance error $R_k = \Sigma - \Sigma_k$:
\[
    R_{k+1} = \Sigma - \Sigma_{k+1} = R_k - B_k R_k B_k.
\]
\end{calculation}

%=============================================================================
\section{Convergence Analysis}
%=============================================================================

\subsection{Contraction of the Covariance Error}

We need to show that $\|R_{k+1}\| \leq c\,\|R_k\|$ for some $c < 1$.

\begin{calculation}[title={Bounding $\|R_{k+1}\| = \|R_k - B_k R_k B_k\|$}]
\textbf{Claim:} For any symmetric matrix $X$ and symmetric $B$ with $0 \prec B \prec I$:
\[
    \|X - BXB\| \leq (1 - \lambda_{\min}(B)^2)\,\|X\|.
\]

\textbf{Proof sketch:} Decompose $X = X_+ + X_-$ where $X_+ \succeq 0$ and $X_- \preceq 0$ (positive and negative parts in the spectral decomposition). Then:
\begin{align}
    X_+ - BX_+B &\succeq 0 \quad\text{and}\quad X_+ - BX_+B \preceq (1-\lambda_{\min}(B)^2)\,X_+, \notag \\
    X_- - BX_-B &\preceq 0 \quad\text{and}\quad X_- - BX_-B \succeq (1-\lambda_{\min}(B)^2)\,X_-. \notag
\end{align}
The first line holds because for any $v$: $v^\top(X_+ - BX_+B)v = v^\top X_+ v - (Bv)^\top X_+(Bv) \leq (1-\lambda_{\min}(B)^2)\,v^\top X_+ v$. Combining gives $\|X-BXB\| \leq (1-\lambda_{\min}(B)^2)\|X\|$.

\medskip
\textbf{Applying this:} We need a lower bound on $\lambda_{\min}(B_k)$.
\end{calculation}

\subsection{Lower Bound on $\lambda_{\min}(B_k)$}

\begin{calculation}[title={Showing $\lambda_{\min}(\Sigma_k) \geq \eta$ for all $k$}]
Define $\eta = \min(\lambda_{\min}(\Sigma), \lambda_{\min}(\Sigma_0))$ and assume $\eta > 0$.

\medskip
\textbf{Proof by induction.} Suppose $\lambda_{\min}(\Sigma_k) \geq \eta$. From the update:
\begin{align}
    \Sigma_{k+1} &= (I - B_k^2)\Sigma_k + B_k\Sigma B_k \notag \\
    &\succeq \lambda_{\min}(\Sigma_k)(I - B_k^2) + \lambda_{\min}(\Sigma)\,B_k^2 \notag \\
    &\succeq \min\bigl(\lambda_{\min}(\Sigma_k),\;\lambda_{\min}(\Sigma)\bigr)\,I \notag \\
    &\geq \eta\, I. \notag
\end{align}
The second line uses $\Sigma_k \succeq \lambda_{\min}(\Sigma_k)\,I$ and $\Sigma \succeq \lambda_{\min}(\Sigma)\,I$. The third line uses $(I-B_k^2) + B_k^2 = I$.

\medskip
\textbf{Therefore:}
\[
    \lambda_{\min}(B_k) = \left(\frac{\lambda_{\min}(\Sigma_k)}{1+\lambda_{\min}(\Sigma_k)}\right)^{(1+\epsilon)/2} \geq \left(\frac{\eta}{1+\eta}\right)^{(1+\epsilon)/2}.
\]
\end{calculation}

\subsection{The Convergence Rate}

Putting it all together:

\begin{proposition}[Proposition 13 in the paper]
Let $\eta = \min(\lambda_{\min}(\Sigma), \lambda_{\min}(\Sigma_0)) > 0$. Then:
\[
    \|\Sigma - \Sigma_k\| \leq \Bigl(1 - \eta^{1+\epsilon}(1+\eta)^{-1-\epsilon}\Bigr)^k\,\|\Sigma - \Sigma_0\|,
\]
\[
    \|b - b_k\|^2 \leq \Bigl(1 - \eta^{1+\epsilon}(1+\eta)^{-1-\epsilon}\Bigr)^k\,\|b - b_0\|^2.
\]
\end{proposition}

\begin{calculation}[title={Deriving the rate}]
The contraction factor per iteration is:
\[
    c = 1 - \lambda_{\min}(B_k)^2 \leq 1 - \left(\frac{\eta}{1+\eta}\right)^{1+\epsilon}.
\]
This is because $\lambda_{\min}(B_k)^2 = \left(\frac{\lambda_{\min}(\Sigma_k)}{1+\lambda_{\min}(\Sigma_k)}\right)^{1+\epsilon} \geq \left(\frac{\eta}{1+\eta}\right)^{1+\epsilon}$.

\medskip
Writing $\left(\frac{\eta}{1+\eta}\right)^{1+\epsilon} = \eta^{1+\epsilon}(1+\eta)^{-(1+\epsilon)}$, we get the rate stated above.
\end{calculation}

%=============================================================================
\section{ODE vs.\ SDE: Why $\epsilon = 0$ Is Optimal}
%=============================================================================

\begin{keyidea}
The convergence rate is $\bigl(1 - \eta^{1+\epsilon}(1+\eta)^{-1-\epsilon}\bigr)^k$. As a function of $\epsilon$, the factor we subtract is:
\[
    \left(\frac{\eta}{1+\eta}\right)^{1+\epsilon}.
\]
Since $\frac{\eta}{1+\eta} < 1$, increasing $\epsilon$ makes this quantity \emph{smaller}, which makes the contraction \emph{worse}. The optimal choice is $\epsilon = 0$ (ODE), giving the contraction factor:
\[
    1 - \frac{\eta}{1+\eta} = \frac{1}{1+\eta}.
\]
\end{keyidea}

\begin{intuition}
Why does the ODE beat the SDE? The SDE injects extra noise during the reverse transport. While this can help with robustness in harder problems (e.g., the two-moons synthetic example with high noise), for the Gaussian case the added stochasticity is \emph{pure overhead} that slows convergence.
\end{intuition}

\subsection{Behavior in Different SNR Regimes}

\begin{itemize}
    \item \textbf{Low SNR ($\eta \ll 1$):} The signal is weak. The ODE rate is $(1-\eta)^k$ while the EM/SDE($\epsilon=1$) rate is $(1-\eta^2)^k$. Since $\eta^2 \ll \eta$, ODE converges \emph{much} faster.

    \smallskip
    \textit{Example:} If $\eta = 0.01$ (very noisy), after $k=100$ iterations:
    \begin{itemize}
        \item ODE error factor: $(1-0.01)^{100} \approx 0.37$
        \item EM error factor: $(1-0.0001)^{100} \approx 0.99$ (barely moved!)
    \end{itemize}

    \item \textbf{High SNR ($\eta \gg 1$):} The contraction is $1/(1+\eta) \ll 1$, so convergence is fast regardless. The corruption is mild, and one iteration nearly recovers $\pi$.
\end{itemize}

%=============================================================================
\section{Comparison with the EM Algorithm}
%=============================================================================

\subsection{What Does EM Do?}

The EM algorithm from [RALL24] updates a current prior $\tilde{\pi}_k = \mathcal{N}(0, \tilde{\Sigma}_k)$ (focusing on centered Gaussians for simplicity) via:

\begin{enumerate}
    \item \textbf{E-step:} Compute the posterior $p_k(x|y)$. For Gaussian prior + Gaussian noise:
    \[
        p_k(x|y) = \mathcal{N}\bigl(\tilde{\Sigma}_k(\tilde{\Sigma}_k+I)^{-1}y,\;\; \tilde{\Sigma}_k - \tilde{\Sigma}_k(\tilde{\Sigma}_k+I)^{-1}\tilde{\Sigma}_k\bigr).
    \]

    \item \textbf{M-step:} The new prior is the marginal of the posterior over the data:
    \[
        \pi_{k+1}(x) = \mathbb{E}_{y\sim\mu}[p_k(x|y)].
    \]
\end{enumerate}

\subsection{The EM Covariance Update}

The paper shows that the EM update for covariances is:
\[
    \tilde{\Sigma}_{k+1} = \tilde{\Sigma}_k + M_k(\Sigma - \tilde{\Sigma}_k)M_k, \qquad M_k = \tilde{\Sigma}_k(\tilde{\Sigma}_k + I)^{-1}.
\]

\begin{keyidea}
Compare with the SCSI update:
\[
    \Sigma_{k+1} = \Sigma_k + B_k(\Sigma - \Sigma_k)B_k, \qquad B_k = \bigl(\Sigma_k(\Sigma_k+I)^{-1}\bigr)^{(1+\epsilon)/2}.
\]
Setting $\epsilon = 1$ gives $B_k = (\Sigma_k(\Sigma_k+I)^{-1})^1 = M_k$. So:
\begin{center}
\fbox{\textbf{EM is exactly the SCSI iteration with $\epsilon = 1$.}}
\end{center}
Since the optimal $\epsilon$ is 0, SCSI with ODE transport is \emph{provably faster} than EM in this setting!
\end{keyidea}

\subsection{Non-Asymptotic Rates}

Beyond the asymptotic (per-iteration) rate, the paper notes:
\begin{itemize}
    \item EM enjoys a non-asymptotic rate of $O(1/k)$ due to its mirror-descent structure.
    \item Empirically, the ODE scheme achieves $O(1/k^2)$, reminiscent of Nesterov-type acceleration in convex optimization.
\end{itemize}
This is what Figure 2 in the paper illustrates: on a log-log plot, EM error decays with slope $-1$ (i.e., $1/k$) while ODE error decays with slope $-2$ (i.e., $1/k^2$).

%=============================================================================
\section{Verifying the Theoretical Assumptions}
%=============================================================================

\subsection{Wasserstein Lipschitz Stability (Assumption 5)}

Assumption 5 requires: $\mathbb{E}\|\Phi_\pi(y) - \Phi_{\tilde{\pi}}(y)\|^2 \leq R\, W_2^2(\pi, \tilde{\pi})$ for $R < 1$.

\medskip
For Gaussians $\mathcal{N}(0,A)$ and $\mathcal{N}(0,B)$ with the ODE ($\epsilon=0$), the transport cost is:
\[
    \mathcal{T}(A;B) := \text{Tr}\bigl((T_A - T_B)(A+I)(T_A - T_B)\bigr),
\]
where $T_A = (A(A+I)^{-1})^{1/2}$ is the ODE solution map.

The Wasserstein-2 distance between $\mathcal{N}(0,A)$ and $\mathcal{N}(0,B)$ is:
\[
    W_2^2(\mathcal{N}(0,A), \mathcal{N}(0,B)) = \text{Tr}(A + B - 2(A^{1/2}BA^{1/2})^{1/2}).
\]

\begin{intuition}
The paper shows numerically (Figure 3) that $\mathcal{T}(A;B) < W_2^2(A,B)$ with overwhelming probability for macroscopic SNR. In the extreme low-SNR limit ($\|A\|, \|B\| \to 0$), one can show $R$ approaches 1 from below, so convergence slows but doesn't fail.
\end{intuition}

\subsection{Condition Number (for KL Guarantees)}

The condition number $\chi$ from Section 4 can be bounded explicitly:
\[
    \chi \leq (1 + \lambda_{\min}(\Sigma)^{-1})^2 \approx \eta^{-2} \quad\text{for small }\eta.
\]
This captures the ill-conditioning of the inverse problem: when $\eta$ is small (low SNR), the condition number is large, and more iterations are needed.

%=============================================================================
\section{Summary}
%=============================================================================

\begin{center}
\renewcommand{\arraystretch}{1.4}
\begin{tabular}{|l|c|c|}
\hline
& \textbf{SCSI (ODE, $\epsilon=0$)} & \textbf{EM ($\equiv$ SCSI with $\epsilon=1$)} \\
\hline
Contraction factor & $\dfrac{1}{1+\eta}$ & $1 - \left(\dfrac{\eta}{1+\eta}\right)^2$ \\[8pt]
\hline
Low-SNR rate & $(1-\eta)^k$ & $(1-\eta^2)^k$ \\
\hline
Non-asymptotic & $O(1/k^2)$ \text{(empirical)} & $O(1/k)$ \\
\hline
\end{tabular}
\end{center}

\medskip

The Gaussian AWGN case study shows that SCSI with ODE transport is the \emph{optimal} choice in this family of methods, converging faster than both the SDE variant and the classical EM algorithm, especially in the challenging low-SNR regime.

\end{document}